\documentclass[11pt,a4paper]{article}
\usepackage[T1]{fontenc}
\usepackage{lmodern,amsmath,amssymb,amsthm,mathtools,hyperref}
\usepackage[margin=27mm]{geometry}
\setlength{\emergencystretch}{2em}
\newtheorem{theorem}{Theorem}[section]
\newtheorem{lemma}[theorem]{Lemma}
\newtheorem{proposition}[theorem]{Proposition}
\newtheorem{corollary}[theorem]{Corollary}
\newcommand{\Q}{\mathbb Q}\newcommand{\Z}{\mathbb Z}
\newcommand{\A}{\mathbb A}
\DeclareMathOperator{\Norm}{Norm}
\title{Zariski density of the prepared pole witnesses}
\author{}\date{}
\begin{document}\maketitle
\begin{abstract}
The simultaneous elliptic realization underlying the pole certificate
gives more than one rational witness.  Its witnesses are Zariski dense
in the two retained pole parameters.  Consequently finitely many
additional nonzero polynomial conditions in those parameters may be
imposed without strengthening the inherited arithmetic input.  We
prove this by independently varying odd multiples of the two realized
non-torsion points.  A finite-algebra corollary permits simultaneous
nonvanishing of all conjugate second-root radicands whenever the pole
input is already a unit in that algebra.  The result does not assert
that the other source guards are simultaneously units.
\end{abstract}

\section{The realization input}

Fix a prepared rational number $b$ in the sense of~\cite{Dependencies},
and write
\[
 B=17b,\qquad V=B(B-b^2)>0,\qquad V=\theta c^2,
\]
where $\theta$ is the positive squarefree representative and
$c\in\Q^\times$.  For $a\ne0$, put
\[
                      E_a:\quad y^2=x(x-a)(x+3a).
\]
The inherited simultaneous-realization input supplies a positive
rational number $n$ and points
\begin{equation}\label{eq:realization}
 Q_1\in E_n(\Q),\qquad Q_2\in E_{n\theta}(\Q),\qquad
 x(Q_1),x(Q_2)\in B\Q^{\times2},
\end{equation}
outside the two-torsion.  Its construction includes a fresh prime
of odd valuation in $n$ outside the support of $\theta$; thus neither
$n$ nor $n\theta$ is a rational square.  These are the same
realized twists used in the source pole-chart lemma~\cite{Source}.

\begin{lemma}\label{lem:torsion}
If $a\in\Q^\times$ is not a square, then
$E_a(\Q)_{\rm tors}=E_a[2](\Q)$.
\end{lemma}
\begin{proof}
There is full rational two-torsion.  Mazur's torsion theorem
\cite{Mazur} therefore leaves the groups
$\Z/2\Z\times\Z/(2m)\Z$, with $1\le m\le4$.
It suffices to exclude points of orders three and four.

The three-division polynomial is
\[
                3x^4+8ax^3-18a^2x^2-9a^4.
\]
A rational root would give a rational zero of
$3z^4+8z^3-18z^2-9$ on putting $z=x/a$.
This polynomial has no root modulo $5$: its five values are
$1,4,1,3,3$.  A rational zero cannot have negative $5$-adic
valuation, since the leading coefficient is a $5$-adic unit.
Thus no rational zero exists.

If a rational two-torsion point with abscissa $e_i$ is a double,
its differences $e_i-e_j$ and $e_i-e_k$ from the other two roots
must be rational squares.  This necessary condition follows
directly from the doubling formula for a cubic with three
rational roots.  At $e_i=0$, the differences $-a,3a$ have
opposite signs.  At $e_i=-3a$, their ratio is $3/4$, not a
rational square.  At $e_i=a$, they are $a,4a$, which cannot
both be squares under the hypothesis.  Hence no point of order
four exists, and the stated torsion group follows.
\end{proof}

In particular, both points in~\eqref{eq:realization} have infinite
order.  For the density theorem it is enough to assume directly
the existence of the two infinite-order points in that display.
This formulation makes the exact arithmetic dependency explicit.

\section{Two independent infinite sets of witnesses}

The squareclass map
\[
 \alpha_a:E_a(\Q)\longrightarrow\Q^\times/\Q^{\times2},
 \quad O\longmapsto1,\quad (0,0)\longmapsto-3a^2,
 \quad (x,y)\longmapsto x\quad(x\ne0)
\]
is a homomorphism.  This follows from the chord-and-tangent law:
the product of the three abscissas cut out by a line with nonzero
intercept is a square.  The lines through $(0,0)$ give its
displayed exceptional value; see also~\cite[Lemma~2.2]{Fixed}.
Thus every odd multiple of $Q_i$ has the same nonzero abscissa
squareclass $B$.  The abscissas of the positive odd multiples
form an infinite set, since an abscissa determines at most two
points and $Q_i$ has infinite order.

Choose arbitrary positive odd multiples and write
\[
 Q_1'=(Bu^2,y_1)\in E_n(\Q),\qquad
 Q_2'=(Bv^2,y_2)\in E_{n\theta}(\Q),
 \qquad u,v\in\Q^\times.
\]
Define
\begin{equation}\label{eq:witness}
 \begin{gathered}
 e=\frac{n}{Bu^2},\qquad
 p=\frac{y_1}{Bu^3},\qquad s=\frac{cv}{u},\qquad
 q=\frac{c^2y_2}{B^2u^2v}.
 \end{gathered}
\end{equation}
Rescaling the first curve by $u^{-1}$ and the second by $c/u$
gives the identities
\[
                 p^2=B(1-e)(1+3e),\qquad
                 Bq^2=(s^2-eV)(s^2+3eV).
\]
None of $p,q,s$ vanishes, because neither multiple is torsion.
Since $n,B>0$ and the positive abscissa of a non-torsion point
on $E_n$ is greater than $n$, one has $0<e<1$.
Consequently the single rational chart is available:
\begin{equation}\label{eq:chart}
 \tau=\frac{p}{e-1},\qquad
 e=\frac{\tau^2-B}{\tau^2+3B},\qquad
 p=-\frac{4B\tau}{\tau^2+3B}.
\end{equation}

\begin{theorem}\label{thm:density}
Assume the realization input~\eqref{eq:realization}.  The set of
pairs $(\tau,s)\in\Q^2$ obtained from rational pole witnesses
with $pqs\ne0$ is Zariski dense in $\A^2$.
In particular, if $F_1,\ldots,F_r\in\overline\Q[T,S]$ are
nonzero polynomials, there is a witness in the single chart
\eqref{eq:chart} with
\[
                  F_j(\tau,s)\ne0\qquad(1\le j\le r).
\]
\end{theorem}
\begin{proof}
Varying the positive odd multiple $Q_1'$ gives infinitely many
values of $u^2$, and therefore infinitely many values of
$e=n/(Bu^2)$.  Formula~\eqref{eq:chart} shows that the corresponding
$\tau$ values also form an infinite set.

Fix one such choice of $Q_1'$ and hence one value of $\tau$.
Independently varying the positive odd multiple $Q_2'$ gives
infinitely many values of $v^2$, and hence infinitely many
values of $s=cv/u$.

Suppose a polynomial $F(T,S)$ vanishes on every pair so obtained.
For each of the infinitely many fixed values $\tau$, the polynomial
$F(\tau,S)$ has infinitely many zeros and is therefore zero.
Its coefficient polynomials in $T$ vanish at infinitely many
values, so each is zero.  Thus $F=0$, proving density.  Applying
this argument to the nonzero product $F_1\cdots F_r$ proves the
last assertion, equally over any algebraic extension of $\Q$.
\end{proof}

\begin{corollary}\label{cor:several}
For any finite list of independently prepared pole inputs, the
product of their rational witness sets is Zariski dense in the
product of their parameter planes.  Any finite collection of
nonzero polynomial conditions involving all those parameters
can therefore be avoided simultaneously.
\end{corollary}
\begin{proof}
Products of Zariski-dense subsets are Zariski dense, by applying
the one-factor coefficient argument successively.
\end{proof}

\section{Nonvanishing in a finite algebra}

Let $A$ be a finite-dimensional commutative $\Q$-algebra of dimension
$d>0$.  Reducedness is not assumed.  Suppose $\mathcal B$ is a
unit of $A$, and let $\mathcal V\in A$.  For indeterminates $T,S$,
put
\[
 \mathcal D(T)=T^2+3\mathcal B,\qquad
 \mathcal e(T)=\frac{T^2-\mathcal B}{T^2+3\mathcal B}
\]
in the localization where $\mathcal D$ is a unit, and define
\[
 \mathcal R(T,S)=
 \frac{(S^2-\mathcal e(T)\mathcal V)
            (S^2+3\mathcal e(T)\mathcal V)}{\mathcal B}.
\]
The algebra norm is the determinant of multiplication on $A$.

\begin{proposition}\label{prop:radicands}
The polynomial $\Norm_A(\mathcal D(T))$ is monic of degree $2d$.
For every rational $t$ where this norm is nonzero,
\[
               \Norm_A(\mathcal R(t,S))\in\Q[S]
\]
has degree $4d$ and nonzero leading coefficient
$\Norm_A(\mathcal B^{-1})$.
\end{proposition}
\begin{proof}
The multiplication matrix of $\mathcal D(T)$ is
$T^2I+3M_{\mathcal B}$, giving the first assertion.
At an admissible $t$, the radicand $\mathcal R(t,S)$ has degree
four in $S$ and leading coefficient $\mathcal B^{-1}$.
Its multiplication matrix therefore has leading coefficient
$M_{\mathcal B^{-1}}$, whose determinant is nonzero.
The determinant has degree $4d$ with the asserted leading term.
\end{proof}

Suppose a rational character $A\to\Q$ sends $\mathcal B$ to
the actual $B$ of a prepared pole input, and sends $\mathcal V$
to its actual $V$.  Theorem~\ref{thm:density} and
Proposition~\ref{prop:radicands} then allow an actual rational
pole witness with both displayed norms nonzero.  Hence its chart
denominator and its second-root radicand are units simultaneously
at every geometric character of $A$.  This remains true for several
poles by Corollary~\ref{cor:several}.

The hypothesis that $\mathcal B$ is already a unit is essential to
this application.  The density theorem proves avoidance of nonzero
specialized polynomial conditions; it cannot remove a condition
that vanishes identically after the other source coordinates have
been fixed.  In particular, it does not establish simultaneous
nonvanishing of all point-chart discriminants or all source
denominators.  Those require a separate argument.

\begin{thebibliography}{9}
\bibitem{Dependencies}
\emph{Arithmetic dependencies of the construction with 29 variables},
project paper, version 8, 8 November 2026.
\bibitem{Source}
\emph{A defining polynomial in twenty-seven variables:
an existence theorem with a uniform degree bound},
project paper, version 12, 9 November 2026, pole-chart lemma.
\bibitem{Fixed}
\emph{No constant pole parameter is complete},
project paper, version 14, 9 November 2026.
\bibitem{Mazur}
B.~Mazur, \emph{Modular curves and the Eisenstein ideal},
Publications Math\'ematiques de l'IH\'ES \textbf{47} (1977), 33--186,
\url{https://doi.org/10.1007/BF02684339}.
\end{thebibliography}
\end{document}
